\section{Empirical Analysis of Earthquake Effects}\label{sec:analysis_earthquake_effects}



\subsection{Main findings}\label{sec:main_findings}

Damages to students' homes varied based on the quality of their homes and the distance of their hometown from the earthquake's asperity. This suggests that we can estimate the causal impact of earthquake damages on student outcomes by using data from a cohort of students with measurable pre-existing vulnerability to the earthquake but whose outcomes were measured before the earthquake struck. A difference-in-differences estimator exploits the differential correlation between vulnerability and outcomes across cohorts to tease out causal impacts.\footnote{In this section, ``vulnerability" refers to overall vulnerability, measured in damage ratios, which considers both construction quality and distance from the asperity.}\par 


Equation \eqref{regression1} presents the regression model I estimate. I consider two achievement outcomes: GPA, which depends on teachers' grades and is in principle observable to classroom peers, and standardized test scores, which are graded centrally and are not directly observed by peers. I use the damage ratios defined in equation \eqref{eq:measure_damage} to measure pre-existing earthquake vulnerability, which reflects actual home damages only for the cohort exposed to the earthquake. The vector $D_{ic}$ of vulnerability variables comprises the damage ratio of student $i$ and the (leave-one-out) mean and standard deviation of damage ratios in $i$'s classroom $c$. The dummy variable $post_i$ takes on value $1$ if a student belongs to the post-earthquake cohort, the one exposed to the earthquake, and $0$ otherwise. The vector $x_i$ of student characteristics includes a lagged achievement measure (the standardized test score in grade 4). The vector $w_{cs}$ of characteristics of classroom $c$ in school $s$ includes the school building's vulnerability.\footnote{ I do not observe the construction materials of the school building, but I observe the shaking intensity in the school's town. To allow for different shaking-resistance levels depending on construction materials, I include as regressors the shaking in the school's town, the shaking interacted with whether the school is public or private (to account for building quality differences across public and private schools),  the shaking interacted with the cohort dummy, the shaking interacted with both the public school and cohort dummies, and the cohort and public school dummies interacted.} The Table notes contain the full list of regressors. \par 





\begin{equation}\label{regression1}
    y_{ics}=\alpha_0 +\alpha_1\cdot w_{cs} + \alpha_2\cdot  x_i + \beta^{'}\cdot  D_{ic}+  post_i\cdot \left[ \gamma + \delta^{'}\cdot D_{ic}\right] +\epsilon_{ics}.
\end{equation}


 For the pre-earthquake cohort, the parameter $\beta$ captures the spurious relationship between vector $D_{ic}$ and outcomes: the location and quality of a student's and her classmates' homes could correlate with unobserved outcome determinants. If such spurious relationship is constant across cohorts, an assumption I assess in section \ref{sec:testing}, the $\delta$ parameters reveal the effects of earthquake damages on achievement, keeping school building damages fixed. This is the parameter vector of interest. \par 

Table \ref{tab:effectsts} presents the results. The outcome in column (1) is the average between the Mathematics and Language SIMCE standardized test scores in eighth grade. Appendix Table \ref{tab:effectsspanmath} shows that considering the two subjects separately yields similar patterns, with somewhat stronger effects on Language at the point estimates. The outcome in column (2) is the GPA in eighth grade, standardized. A one standard deviation increase in a student's damage ratio, corresponding to increasing the collapsed portion of the home by 4.2 percentage points (around USD $3,500$ in damages), lowers test scores by 0.03 standard deviations (std) and GPA by 0.02 std, albeit insignificantly for GPA. To put the magnitude into perspective, this impact is a fifth of that of a one-standard-deviation improvement in teacher value added (\cite{chetty2014measuring}).\footnote{Teachers are one of the most important school inputs into the production of achievement, but school inputs are generally not as impactful as home interventions (e.g. \cite*{zhou2022impacts}, \cite{heckman2006skill}).}\par 

\input{output/tables/main_effects_ts_GPA}


The average damages to the homes of classmates have positive effects on own test scores and GPA. An increase in mean peer damages of one standard deviation of the damage distribution (i.e., a 4.2 percentage point increase in the portion of the home that collapsed) increased test scores by 0.05 standard deviations, and GPA by 0.04 standard deviations. This suggests that schools counteracted any potential adverse learning conditions caused by average damages. Overcompensation in response to the earthquake was documented also in post-earthquake crime prevention in Chilean municipalities (\cite{hombrados2020lasting}). Classroom-level damage dispersion had negative effects on test scores and GPA, of similar magnitudes. An increase in the within-classroom standard deviation of damages of one standard deviation of the damage distribution (i.e., a 4.2 percentage point increase in the portion of the home that collapsed) lowered test scores and GPA by around 0.085 standard deviations. These results suggest that schools did not entirely compensate detrimental effects on student learning due to damage dispersion within classrooms. \par 


To summarize, damages affected the learning of the student living in the damaged home. The detrimental impacts occurred at a critical time in the educational path of students (the year before transferring to secondary education) and were disproportionately borne by students of lower socioeconomic status due to their greater exposure (Figures \ref{fig:vc_by_peduc} and \ref{fig:damage_propagation_by_SES}). While schools could not mitigate the impact of such individual-level shocks, they appear to have successfully mitigated the effect of the average level of damage in the classrooms. By contrast, dispersion in damages across classmates had negative effects on learning outcomes, on average.\par 


\subsection{Identifying assumption}\label{sec:testing}


  The identifying assumption underlying the estimator is that the relationship between achievement and the earthquake vulnerability variables would be the same in the pre- and post-earthquake cohorts in the absence of the earthquake. A concern is that the estimates may capture changes in this relationship across cohorts, rather than true damage impacts. For example, the estimate of the impact of mean damages in the classroom would be biased if the government introduced a policy between 2009 and 2011, the period between the outcomes for the two cohorts were measured, that changed the student composition across schools, such as changes to the vouchers for disadvantaged students to attend private schools.\footnote{Chile has a voucher policy in place, but it did not undergo any changes at this time (\cite{neilson2021targeted}).} If such a policy were introduced, it could alter how a school's mean damage measure, based on the socioeconomic composition, correlates with its unobserved quality across cohorts. This would violate the identifying assumption. To address such concerns, all specifications include a set of controls for socioeconomic composition. By the same logic, they include controls for individual characteristics.\footnote{Appendix Table \ref{tab:effectstscontrolsyn} shows that effect estimates are slightly larger but broadly similar in specifications without such individual and group controls, retaining only the three individual characteristics used to build the damage measure.} Identification, therefore, relies on variation across students and classrooms in earthquake exposure, keeping fixed student characteristics and classroom characteristics such as socioeconomic status.\par

As the earthquake occurred a few days before the start of grade 7 and outcomes are measured in grade 8, a second concern is that the estimates may capture the effect of a reallocation of students across classrooms and/or schools between grades 7 and 8, occurring in response to earthquake damages, that the strategy above and the exclusion from the sample of forced relocations immediately after the earthquake (Section \ref{sec:data}) fail to account for. To examine this, I tracked the movement of students across schools and classrooms between the $7^{th}$ and $8^{th}$ grades, for both cohorts.\par

Table \ref{tab:descriptiveswitches} reports the descriptive analysis. Switches between the $7^{th}$ and $8^{th}$ grades are not common: only $11\%$ of students in the whole sample changed either school or classroom between these grades. This is aligned with what we would expect: students typically move school one grade later, during the transition from primary to secondary school. The vast majority of switches concern school changes: $8.6\%$ of students changed school, while conditional on staying in the same school, only $2.8\%$ changed classroom. Focusing on students and schools in earthquake regions and residing more than 1 km from the coast, the main sample of analysis, I find similar frequencies of all types of switches, as can be seen in the right panel of Table  \ref{tab:descriptiveswitches}. \par 


Next, I analyze whether students could have moved between the $7^{th}$ and $8^{th}$ grade in response to the earthquake. In the main analysis sample (earthquake regions, non-coastal towns), the fraction of overall switches, school switches, and classroom-within-school switches is identical in the pre- and post-earthquake cohorts, as shown in Panel A of Table \ref{tab:effectsswitch}. This suggests that changes to schools or classroom enrollments between these grades were not a margin of response to the earthquake.\par


\input{output/tables/effects_on_switches}


 To complement the before-after analysis, I implemented a difference-in-differences analysis that examines whether the change in switches across cohorts differed between the regions afflicted and those not afflicted by the earthquake. This additional analysis delivers precise zero effects as well, as can be seen in Panel B of Table \ref{tab:effectsswitch}, where only the coefficient in column (3) is significant ($p<0.10$), but it is small and the effect becomes a precise zero once controls are included in column (4). The evidence, therefore, suggests that the few reallocations observed in the post-earthquake cohort and in regions affected by the earthquake were not induced by the earthquake itself, but rather fell within the typical reallocations we observe in normal years around the country.\par 


Next, I examine the effects of the earthquake-damage measures used in the main analysis on lagged outcome measures, as a pre-trend test and as a way to further assess compositional impacts. This test shows precise zero effects of the damage to a student's own home, of the average damages in the classroom, and of the standard deviation of damages in the classroom on lagged (fourth-grade) test scores and GPA, as shown in Table  \ref{tab:effectslaggedtsGPA}. These results give us further confidence that the findings are not driven by confounding effects due to changes in observed student composition.\par 



\input{output/tables/main_effects_ts_2025_lagged}

Bias could still arise if unmeasured components of socioeconomic composition or of individual characteristics correlate with earthquake vulnerability variables differently across cohorts. To address this concern, I assess the validity of the identifying assumption using data from regions of Chile not affected by the earthquake. As Table \ref{summaryprepost} showed, around a quarter of students in the sample lived in regions so far from the asperity that no damage to buildings occurred. \par 

The ideal test would be to re-estimate equation \eqref{regression1} on the sample of students living in non-affected regions. Failing to reject the null hypotheses that $\delta=0$ provides evidence in favor of the identifying assumption: we would not be able to reject the notion that, in the absence of an earthquake, the measures of earthquake vulnerability used in the analysis (own vulnerability, and mean and standard deviation of vulnerability among classmates) correlate with outcomes identically across cohorts (under the assumption that the evolution of such correlation across cohorts in the regions unaffected by the earthquake equals that in the regions affected).\footnote{More formally, letting $y_{0}$ denote the potential outcome in the absence of the earthquake, and $E=1$ for regions affected by the earthquake and $E=0$ for regions not affected, the identifying assumption is $\nabla_{D} E[y_0|post,E=1]-\nabla_{D}E[y_0|pre,E=1]=\nabla_{D}E[y_0| post,E=0]-\nabla_{D}E[y_0|pre,E=0]$, $\forall D$, where $\nabla_{D}$ denotes the gradient with respect to the components of $D$, the vector collecting own seismic vulnerability and the classroom mean and standard deviation of seismic vulnerabilities. }\par 

I cannot run the ideal test because damage ratios are equal to zero by construction in regions where the ground did not shake ($p^m(vc=j,I_i=0)$ in equation \eqref{eq:measure_damage} is zero for all $j,m$). As a result, I focus on variation in students' home quality, which can be constructed for any student nationwide. When holding the town of residence constant, this metric becomes a proxy for earthquake vulnerability. This is because within a town, differences in damage ratios are determined solely by differences in housing quality.\par 
Therefore, I use the sample of classrooms where every student resides in the same town (the school's town), and include dummy variables for students' town of residence, uninteracted and interacted with the cohort dummy. I then re-estimate regression \eqref{regression1}, using earthquake vulnerability measures based on housing quality in vector $D_{ic}$. For each student, we have a vector of probabilities indicating the likelihood their home falls into one of three seismic vulnerability classes. From this vector I construct an index. A value of 1 indicates that a student certainly lives in a high-vulnerability home, a value of 0 that she certainly lives in a low-vulnerability home.\footnote{The index is $1 \cdot \hat{p}_i^{HV} + 0.5 \cdot  \hat{p}_i^{MV} + 0 \cdot \hat{p}_i^{LV}$.} I standardize this index across the entire sample, so that a one-unit increase corresponds to an increase in earthquake vulnerability by one standard deviation. I also generate the leave-one-out classroom mean and standard deviation of this vulnerability index. \par 

Table \ref{tab:testidentassmpt} shows the results. As a plausibility check on the measure of earthquake vulnerability only based on housing quality, the first two columns are based on the sample of students from earthquake-affected regions. The patterns align with the main findings presented in Table \ref{tab:effectsts}, suggesting that the measure of damages based on housing quality and keeping location fixed is a good proxy for the measure based on damage ratios used in the main analysis.\footnote{The main difference is the significant positive impact of the standard deviation. This can be explained by the fact that damage dispersion has heterogeneous effects by baseline test scores (positive for low-performing and negative for high-performing students, as per section \ref{sec:het_ability}), and the sample underlying  Table \ref{tab:testidentassmpt} is selected (those attending schools that do not attract students from other towns tend to be lower-performing students).}\par 


\input{output/tables/test_identif_assmpt}

Columns (3) and (4) assess the validity of the identifying assumption, presenting estimates of $\delta$ from equation \eqref{regression1} using data from regions {\itshape not} affected by the earthquake. In almost all cases, we cannot reject the hypotheses that $\delta=0$ at any conventional significance level, with the exception of a significant effect of own home vulnerability on GPA, but this does not hold for test scores, where the effect of own home vulnerability is smaller and statistically undetectable. This suggests that any potential spurious correlation between pre-existing vulnerability to the earthquake and achievement --- attributable to correlations between housing quality and unobserved achievement determinants --- is constant across cohorts. Importantly, the evidence is consistent with a lack of spurious correlation for the treatment variables of central interest in this study, that is, the classroom-level earthquake damages. This gives us further confidence in interpreting the findings on the spillover effects from peer damages as causal.\par


\subsection{Heterogeneity by baseline test scores }\label{sec:het_ability}

The impacts of damages to students' homes can vary depending on the students' prior achievement. To explore this, I estimate model \eqref{regression1het},  where prior achievement \(a_i\) is measured through a standardized test in fourth grade:
\vspace{-1em}
\begin{align}
y_{ics} &= \alpha_0 +\alpha_1\cdot w_{cs} + \alpha_2\cdot  x_i + \beta_1^{'}\cdot  D_{ic} \nonumber \\
&\quad + \beta_2^{'}\cdot  D_{ic}\cdot a_i+  post_i\cdot \left[ \gamma_1 + \gamma_2 \cdot a_i + \delta_1^{'}\cdot D_{ic} + \delta_2^{'}\cdot D_{ic}\cdot a_{i}\right] +\epsilon_{ics}, \label{regression1het} 
\end{align}

\noindent and where, like before, $x_i$ includes prior achievement.\footnote{As in the previous model, prior achievement enters the specification additively, so $x_i$ includes $a_i$; equation \eqref{regression1het} displays explicitly only the interaction terms involving $a_i$.\label{fn:intera}} The parameters of interest are $\delta_1$, which captures the effects of $D_{ic}$ for a student with mean preparation (i.e., $a_i=0$), and $\delta_2$, the coefficient on the interaction term. These inform us about the variation in the effects across students with different prior achievement. I chose a centrally graded standardized test as the measure of baseline achievement, as such tests are less prone to teacher bias in grading (\cite{carlana2019implicit}) and therefore could be considered more objective than baseline GPA. \par

 The results are presented in Table \ref{tab:heteffectsts} and Figure \ref{fig:me_nofe}. The detrimental impacts of damages to a student's own home did not significantly vary with a student's baseline achievement, as seen in the second row of Table \ref{tab:heteffectsts}. Similarly, the effects  of the average damage among classmates varied insignificantly with baseline achievement, as seen in the fourth row. Average damage had insignificantly stronger positive impacts on the test scores of higher-baseline-achievement students (first column), for whom the impacts on test scores of mean damages were positive and statistically significant, as seen in the top-left panel of Figure \ref{fig:me_nofe}. This suggests that any remedial measures undertaken by schools may have benefited the test scores of higher-baseline-achievement students more, although differences between students are imprecisely estimated. There is no heterogeneity on the impacts of average damages on GPA. \par 
\input{output/tables/het_effects_ts_GPA_bysimce.tex}


\begin{figure}[h!] 
 	\centering
 	\includegraphics[scale=0.18]{output/figures/heteffects_ts_GPA_by_simce_nofe.png} %0.25
 	\caption{Marginal effects on standardized eighth-grade test score and GPA by baseline test score. {\itshape Notes}: Marginal effects of leave-one-out average damage among classmates and leave-one-out standard deviation of damage among classmates. Effects obtained from estimating the regression model in equation \eqref{regression1het}. 80\% and 90\% confidence intervals reported. }
 	\label{fig:me_nofe}
 \end{figure}





While the dispersion in damages in the classroom showed negative effects on average (Table \ref{tab:effectsts}), the effects varied substantially and significantly across students (last row of Table \ref{tab:heteffectsts}). A rise in such dispersion raised the achievement of students with lower baseline achievement and lowered that of students with higher baseline achievement, as can be seen in the second column of Figure \ref{fig:me_nofe}. These results hold regardless of the outcome measure, eight grade test score or GPA. For some students, the dispersion in damages had a similar or even larger effect than that of the damages at their own home.\footnote{Magnitude comparisons rely on the unit of measurement. In Table \ref{tab:heteffectsts}, all treatment variables are expressed in standard deviations of the overall damage distribution. The table reports the impacts of increasing each treatment variable by the same amount in absolute terms, that is, increasing the portion of the home that collapsed by around 4.4 percentage points, or increasing the reconstruction costs by around USD 3,600. In Tables \ref{tab:effectstsst} and \ref{tab:heteffectstsst}, each treatment variable is instead standardized by its own distribution. The conclusion that some students suffered similar or larger impacts from the damage dispersion in the classroom than from the damage to their own home holds in both cases. \label{fn:magnitudecompar} }\par 


These findings hold regardless of the specification, interaction terms, and baseline achievement measure used. Relaxing the linearity assumption using an interaction with deciles of the baseline test score distribution results in less precise estimates but confirms the patterns (Table \ref{tab:heteffectsdecilesall}).\footnote{In principle, more flexible non-parametric approaches could be used to model the bias arising from unobserved correlates within the difference-in-differences framework, as demonstrated in the seminal conditional difference-in-differences method developed in \cite*{heckman1998characterizing}. In the context of social effects, this could be achieved by relaxing parametric restrictions of control function approaches (see \cite{brock2001interactions,durlauf2006multinomial}, who, by bringing the insights from \cite{heckman1979sample} and \cite{heckman1986alternative} into the study of social effects, demonstrated that control functions can aide in their identification.). The treatment effects could be modelled as non-parametric functions of student characteristics to examine heterogeneity more flexibly. However, such non-parametric methods deliver impractically large estimator variances in this empirical setting. } Including interactions with other student socioeconomic characteristics does not change the findings (Table \ref{tab:heteffectstsallint}). Replacing the fourth grade test score with fourth grade GPA to measure baseline achievement yields broadly similar results (Appendix Figure \ref{fig:me_nofe_byGPA4}).   \par



In summary, the negative effects of damages to a student's own home were similar across the baseline achievement distribution. The positive effects of classroom mean damages were insignificantly stronger for the test scores of higher-performing students. Relatively substantial earthquake impacts arose from damage dispersion in classrooms, especially lowering the achievement of students with high prior achievement.\par 





\subsection{Robustness}\label{sec:robustness}

The analyses restricted the sample to non-coastal towns to mitigate potential attenuation bias from damages from the tsunami, which are not adequately accounted for by damage ratios. A town is defined as coastal if it is within a 1 km strip of the coast. I repeated the analyses defining coastal proximity as within 0.5 and 1.5 km of the coast. I also repeated the analyses in the unrestricted sample that includes coastal towns. As shown in Appendix Table \ref{tab:effectststsunami}, the results are robust to different definitions of coastal proximity. The conclusions stand even considering the entire sample, but, as expected, whenever damages enter linearly, estimates are attenuated towards zero in this case.\par 



The analyses allowed for correlation in the error terms of an unknown form between students in the same school and cohort. But error terms of students in different schools that are geographically close may correlate, as shaking is similar in nearby schools. I employ two methods to account for this. First, I cluster the standard errors at the school-town-by-cohort rather than the school-by-cohort level. Second, I use Conley's method (\cite{conley1999gmm}) to compute the standard errors, allowing for a more flexible spatial correlation in the residuals. The method assumes that the spatial dependence between two students residing in different towns is a decreasing function of the distance between the towns, and that beyond a pre-specified distance cutoff, there is no dependence. I present results for different cutoff distances (10 km, 25 km, 50 km, 250 km). As shown in Appendix Tables \ref{tab:effectstsspatcor} and \ref{tab:effectstsconley}, the standard errors are similar across methods and distance thresholds, suggesting that clustering at the school-by-cohort level accurately captures the spatial correlation in the data-generating process. 






